\documentclass[10pt]{article}
\usepackage[a4paper,margin=0.9in]{geometry}
\usepackage{amsmath,amssymb}
\usepackage[hidelinks]{hyperref}
\usepackage{parskip}
\setlength{\parskip}{0.55em}
\pagestyle{empty}

\begin{document}

\begin{center}
{\large\bfseries SpectraForge: An Open, Physics-Grounded Benchmark for Band Selection\\
in Multi-Excitation (4D) Fluorescence Hyperspectral Imaging}\\[0.6em]
{\normalsize Narek Meloyan\textsuperscript{1,2,3}, Narine Sarvazyan\textsuperscript{1,3,4}}\\[0.2em]
{\footnotesize \textsuperscript{1}Zaven P.\ and Sonia Akian College of Science and Engineering, American University of Armenia, Yerevan, Armenia\\
\textsuperscript{2}Engineering Research Center, American University of Armenia \quad
\textsuperscript{3}L.A.\ Orbeli Institute of Physiology NAS RA, Yerevan, Armenia \quad
\textsuperscript{4}George Washington University, Washington, DC, USA\\
\texttt{narek\_meloyan@edu.aua.am}, \texttt{phynas@gwu.edu}}
\end{center}

\vspace{0.3em}

Multi-excitation fluorescence hyperspectral imaging records, for every pixel, a
four-dimensional excitation-emission cube, and choosing which spectral bands to keep is
central to cheaper acquisition and compact classification. Progress on band selection for
this modality is held back by a data problem: labeled real 4D fluorescence cubes are
scarce, and even when available they carry no ground truth, so a selector cannot be scored
against the bands it should have found. We present \textbf{SpectraForge}, an open,
physics-grounded generator that renders 4D excitation-emission cubes from a parametric
fluorophore model obeying the standard laws of fluorescence photophysics (Kasha, Vavilov,
Franck-Condon, Stokes shift), with optional inner-filter attenuation, reabsorption, FRET,
scatter, photon and read noise, and fixed-pattern clutter. Every cube carries perfect
ground truth (concentration maps and true emission peaks), parametric dyes can be replaced
by real measured spectra imported from FPbase, and an interactive workbench lets a user
paint scenes, render, and validate a selection against ground truth. Crucially, scenes are
built so that band variance is decoupled from class informativeness (measured correlation
$\approx -0.1$), which makes the benchmark an honest test rather than a trivial one, at
difficulty controllable from a single call.

Alongside the generator we define an evaluation protocol whose two decisive controls, a
random-selection baseline and the full spectrum, separate genuine gains from
self-deception, together with a classifier panel and a nonlinear-only information metric
for regimes where linear scoring cannot credit a nonlinear-shape selection.

Using SpectraForge we characterize autoencoder-plus-latent-perturbation band selection.
The benchmark exposed a failure invisible to raw loss (a published model that selected at
chance, macro-F1 $\approx 0.33$) and traced it to four trainer defects and an emission-axis
band collapse; correcting these recovers the selector to $85$--$97\%$ of a labels-using
oracle and matches the strongest classical blind selector (PCA loadings). A
significance-tested battery then shows the learned selector \emph{ties} PCA loadings on
realistic and FRET data and \emph{significantly exceeds} it in the nonlinear reabsorption
regime (macro-F1 $0.449$ vs $0.431$, $95\%$ CI $[+0.001,+0.029]$), capturing nonlinear-only
information that linear selection misses. The same controls delimit where unsupervised
selection stops helping, and reveal that a convolutional model is structurally necessary
when discrimination is carried by spatial texture, where per-pixel selection is provably
blind (below random) while a spatial model reaches the oracle.

SpectraForge turns ``does this selector help'' into a measured, reproducible question for
4D fluorescence imaging. The generator, the evaluation harness, and the full experiment
record are released, and validation on real Lichens and Collagen cubes with the same
controls is the immediate next step.

\end{document}
